\documentclass[12pt]{article}
\usepackage[utf8]{inputenc}
\usepackage[T1]{fontenc}
\usepackage{hyperref}
\title{Panobo Language Resources for K-12 Curriculum Development}
\author{Compiled by Claude Sonnet 5.5 (claude-sonnet-5-5) for Andr\'as Kornai}
\date{}
\begin{document}
\maketitle

\section{Introduction}
Panobo (ISO 639-3 pno) is an extinct or nearly extinct Panoan language of Peru, in sources also called Wariapano, Pano, Panavarro. It was spoken in Ucayali department. Glottolog lists it as extinct (citing Moseley 2010); Wikipedia reports the last speaker of the Pano dialect died in 1991 while also giving a ``moribund (2013)'' status. Search found no educational or revitalization materials.

\section{Primary Educational Resources}
None found in search results.

\section{Linguistic and Reference Materials}
\begin{itemize}
\item Wariapano (Panobo) language, Wikipedia: \url{https://en.wikipedia.org/wiki/Wariapano_language}
\item Panoan languages, Wikipedia: \url{https://en.wikipedia.org/wiki/Panoan_languages}
\item Glottolog entry: \url{https://glottolog.org/resource/languoid/id/pano1255}
\item Numeral system entry for Panobo (Max Planck EVA; word list source attributed to SIL's Stephen Parker): \url{https://lingweb.eva.mpg.de/channumerals/Panobo.htm}
\end{itemize}

\section{Religious and Community Resources}
None found.

\section{Digital and Audio Resources}
None found.

\section{Curriculum Development Context}
Documentation is limited to word lists and reference entries. A heritage curriculum would have to rely on related living Panoan languages and on archival materials, which this search did not locate.
\end{document}